\chapter{Identidad PGP e intercambio de información}
\section{Diseño y ciclo de vida de la clave}
\pendiente{Explicar la relación entre OpenPGP y la implementación elegida. Justificar identidad, parámetros, caducidad, copia protegida y revocación con documentación vigente. Crear una clave propia y registrar su huella pública.}
\section{Publicación y verificación de la clave pública}
\pendiente{Elegir un servidor público compatible, publicar solo la clave pública y demostrar su consulta o recuperación. Documentar fecha y huella. No confundir publicación de una clave con verificación de la identidad de su titular.}
\section{Intercambio con compañeros}
\pendiente{Obtener y verificar por un canal adecuado las huellas públicas de los compañeros. Documentar un envío cifrado y una recepción que se descifra correctamente. Diferenciar la prueba de envío de la confirmación efectiva de recepción.}
\section{Firma y análisis de resultados}
\pendiente{Si se añade firma como mejora, demostrar su verificación y explicar qué aporta frente al cifrado. Registrar incidencias, límites y aprendizaje. La firma adicional es una propuesta, no un requisito explícito del enunciado.}
